\chapter{Conclusiones}
\label{chap:conclusiones}

\section{Conclusiones principales}

El trabajo ha construido una memoria tecnica y academica para un sistema de estimacion explicable de volatilidad implicita en opciones sobre el IBEX. La cadena parte de operaciones y contratos, asocia futuros subyacentes, calcula tipos compuestos, invierte Black-76 y produce una base final con 312.615 observaciones validas. Sobre esa base se entrenan familias de modelos y se genera una capa de explicabilidad orientada a uso operativo.

La primera conclusion es que la calidad del dato es el nucleo del problema. La volatilidad implicita no aparece en las fuentes originales; se construye. Por tanto, los filtros de precio, cantidad, vencimiento, strike, subyacente temporalmente valido y solver numerico son tan importantes como el algoritmo de aprendizaje.

La segunda conclusion es que la no linealidad importa. La regresion lineal ofrece una referencia interpretable, pero su rendimiento fuera de muestra queda lejos de Random Forest y XGBoost. Esto confirma que la superficie de volatilidad requiere capturar interacciones y curvaturas.

La tercera conclusion es que el mejor modelo cuantitativo disponible fue Random Forest, con RMSE de test 0,088403 y $R^2$ 0,790941. XGBoost obtuvo un resultado cercano. La red tensor-train progresiva mostro una mejora notable frente a su version no progresiva, pero no alcanzo a los ensembles de arboles.

La cuarta conclusion es que la explicabilidad debe evaluarse como un sistema, no como un unico grafico. SHAP aporta atribucion, ICE y ALE respuesta funcional, superficies lectura financiera, vecinos soporte local, y surrogates fidelidad aproximada. Cada herramienta cubre una pregunta distinta y sus limitaciones se compensan parcialmente entre si.

\section{Implicaciones practicas}

Desde el punto de vista practico, el proyecto puede servir como base para una herramienta de analisis de volatilidad. Un usuario podria cargar un modelo final, seleccionar una muestra de test o introducir variables manuales, obtener una prediccion y revisar explicaciones locales y globales. La API y el almacenamiento local/GCP apuntan a una arquitectura desplegable.

Para equipos de model risk o validacion, el valor esta en la trazabilidad. La memoria y los anexos identifican de donde sale cada dato, que validaciones se aplican, como se separan las muestras, que metricas se usan y que artefactos explicables se generan. Esta trazabilidad reduce la distancia entre experimento academico y herramienta auditable.

\section{Lineas futuras}

Una primera extension natural es enriquecer las variables con liquidez, volumen, hora intradia, bid-ask spread, volatilidad realizada y medidas de regimen de mercado. Estas variables podrian explicar parte de los residuos extremos.

Una segunda linea es imponer restricciones financieras o regularizaciones sobre la superficie. Por ejemplo, penalizar irregularidades excesivas o incoherencias conocidas podria mejorar robustez, especialmente en zonas poco pobladas.

Una tercera linea es comparar la inversion Black-76 con alternativas o ajustes especificos para dividendos, tipos y convenciones de mercado. La variable objetivo podria cambiar si se modifica la forma de calcular el precio teorico.

Una cuarta linea es ampliar la validacion temporal con backtesting por ventanas moviles y evaluacion por regimenes. El periodo 2017--2022 incluye cambios de tipos y episodios de estres; separarlos podria revelar estabilidad o fragilidad de cada familia.

Finalmente, la explicabilidad podria evaluarse con usuarios. Un dashboard tecnicamente completo no garantiza que las explicaciones sean comprensibles o utiles para distintos perfiles. Diseñar pruebas de uso con analistas, validadores y docentes permitiria priorizar visualizaciones y mejorar comunicacion.

\section{Balance final}

El balance final es positivo con una cautela metodologica. El proyecto demuestra que una arquitectura de datos, modelos y explicabilidad puede aplicarse a volatilidad implicita de opciones IBEX con resultados cuantitativos razonables. Tambien demuestra que la documentacion tecnica existente puede transformarse en una memoria academica coherente, siempre que se separe el relato principal de los anexos operativos.

La cautela es que el modelo aprende una variable construida. La volatilidad implicita depende de Black-76, del subyacente asociado, del tipo compuesto, de la hora de ejecucion y del filtro de operaciones. Cualquier conclusion sobre prediccion debe leerse dentro de ese marco. La fortaleza del trabajo es precisamente que ese marco queda documentado.

\section{Recomendaciones de uso}

Para uso academico, se recomienda emplear la memoria como documento principal y los anexos como guia tecnica. Para uso de desarrollo, se recomienda mantener sincronizada la matriz de trazabilidad con \path{doc/docs/}: si se anade una pagina de documentacion, debe reflejarse en el anexo correspondiente. Para uso de validacion, se recomienda conservar metadatos de cada ejecucion y no sobrescribir resultados sin versionado.

Para uso financiero, se recomienda tratar el dashboard como herramienta de inspeccion. Una prediccion debe acompanarse de su explicacion local, vecinos cercanos y diagnostico de region. Si una muestra cae fuera de soporte, la salida debe interpretarse con prudencia aunque el modelo devuelva un valor numerico.

\section{Trabajo pendiente para una version final entregable}

Antes de entregar una version definitiva, convendria reemplazar los tutores pendientes en \path{variables.sty}, revisar estilo institucional de portada, anadir figuras propias del dashboard si se desea una memoria mas visual y compilar con una distribucion LaTeX instalada. Tambien seria conveniente revisar terminologia con el tutor, especialmente en torno a unidades de tipos y convenciones de mercado.

Estas tareas no cambian la sustancia tecnica del trabajo, pero mejoran presentacion y ajuste formal. La memoria ya recoge la estructura exigida por la guia: titulo, resumen, introduccion, marco teorico, metodologia, resultados y discusion, conclusiones y bibliografia.
